\scanpage{225}
\margin{Feuille volante \unsure{paumée}...}
\uline{Prop 6} Soit $B$ une catégorie $k$-lin. \unsure{karoubienne}.
Alors le foncteur
\[
ND_\bullet:\unsure{B^*_*|B^*_k}\mathbin{-}B_\bullet
\]
est donné (à isom can. près) par
\[
ND_i(L_*)\simeq\operatorname{Hom}_{k,*}
(\underbrace{\Lambda^i\Phi_*}_{\Phi^i_{\unsure{\bullet|*}}},L_*)
\]
(pour la déf. de $\operatorname{Hom}_{k,*}$, cf prop. 3, 5),
\deleted{l'\unsure{opér. différentiel étant}}
\[
ND_\bullet(L_*)\simeq\operatorname{Hom}_{k,*}
(\underbrace{\Lambda^\bullet\Phi_*}_{\Phi_*^\bullet},L_*)
\]
l'opérateur diff. sur $ND_\bullet(L_*)$ provient de l'opér. différentiel sur
$T_{\unsure{\lambda\lambda|\wedge\lambda}}$ sur $\Lambda^\bullet\Phi_*$.

En effet, \deleted{\unsure{par ex.}} soit $l_\bullet=ND_\bullet(L_*)$,
et posons $\Phi(i)_\bullet=ND_\bullet(\deleted{\unsure{\Phi}}\Lambda^i\Phi_*)$,
alors on a par \inserted{l'\unsure{équiv.} de Dold-Puppe}
\[
\operatorname{Hom}_{k,*}(\Lambda^i\Phi_*,L_*)
\simeq\operatorname{Hom}_{k,\bullet}(\Phi(i)_\bullet,l_\bullet).
\]
Or $\Phi(i)_\bullet=C_\bullet^{i-1}$ pour $i\geq1$ (prop 4), et
\inserted{on $\Phi(0)_\bullet=k$ placé en degré $0$}
a bien
\[
\left\{\begin{aligned}
\operatorname{Hom}_{k,\bullet}(C_\bullet^{i-1},l_\bullet)&\simeq l_i\quad(i\geq1)\\
\operatorname{Hom}_{k,\bullet}(\deleted{\unsure{C}}k[0]_\bullet,l_\bullet)&\simeq l_0
\end{aligned}\right.
\]
et il reste à faire de petites vérifications de compatibilité --.

\uline{Cor} Le foncteur $ND^\bullet:B^*\mathbin{-}B^\bullet$ \unsure{s'explicite} comme
\[
ND^\bullet(L^*)\simeq\Lambda^\bullet\Phi_*\otimes_{k,*}L^*
\]
[où pour un $M_*\in\mathrm{Ab}_{k,*}$ et $L^*\in B^*$,
$M_*\otimes_{k,*}L^*$ désigne l'objet de $\operatorname{Hom}(B,(\mathrm{Ab}_k))$
\[
b\longmapsto\operatorname{Hom}_{\mathrm{Ab}_{k,*}}
(M_*,\operatorname{Hom}_{B}(L^*,b))
\]
]

Il suffit d'appliquer la prop. 6 à $B^o$.
\margin{224}
